\documentclass[]{article}
\usepackage[margin=1in]{geometry}
\usepackage{amsmath, amssymb}
\usepackage[most]{tcolorbox}
\usepackage{xcolor}
\usepackage{enumitem}
\usepackage{fancyhdr}

% Page setup
\pagestyle{fancy}
\fancyhf{}
\rhead{AMC12B 2017 Problem 24 Analysis}
\lhead{\today}

% Color definitions
\definecolor{problemconfig}{RGB}{230, 242, 255}
\definecolor{strategic}{RGB}{255, 230, 242}
\definecolor{tactical}{RGB}{242, 255, 230}
\definecolor{techniques}{RGB}{255, 242, 230}
\definecolor{verification}{RGB}{255, 230, 230}
\definecolor{solution}{RGB}{230, 255, 255}
\definecolor{competition}{RGB}{242, 242, 255}
\definecolor{gold}{RGB}{218, 165, 32}

% Box styles
\newtcolorbox{problemconfigbox}{
    colback=problemconfig,
    colframe=blue,
    title=Problem Configuration Analysis,
    fonttitle=\bfseries,
    arc=3pt,
    boxrule=1pt,
    breakable
}

\newtcolorbox{strategicbox}{
    colback=strategic,
    colframe=purple,
    title=Strategic Framework Application,
    fonttitle=\bfseries,
    arc=3pt,
    boxrule=1pt,
    breakable
}

\newtcolorbox{tacticalbox}{
    colback=tactical,
    colframe=green,
    title=Tactical Methodology Selection,
    fonttitle=\bfseries,
    arc=3pt,
    boxrule=1pt,
    breakable
}

\newtcolorbox{techniquesbox}{
    colback=techniques,
    colframe=orange,
    title=Conversion Techniques Application,
    fonttitle=\bfseries,
    arc=3pt,
    boxrule=1pt,
    breakable
}

\newtcolorbox{verificationbox}{
    colback=verification,
    colframe=red,
    title=Verification Strategy Design,
    fonttitle=\bfseries,
    arc=3pt,
    boxrule=1pt,
    breakable
}

\newtcolorbox{solutionbox}{
    colback=solution,
    colframe=cyan,
    title=Solution Roadmap,
    fonttitle=\bfseries,
    arc=3pt,
    boxrule=1pt,
    breakable
}

\newtcolorbox{competitionbox}{
    colback=competition,
    colframe=purple,
    title=Competition Strategy Integration,
    fonttitle=\bfseries,
    arc=3pt,
    boxrule=1pt,
    breakable
}

\newtcolorbox{keyinsightbox}{
    colback=yellow!20,
    colframe=gold,
    title=Key Insight,
    fonttitle=\bfseries,
    arc=3pt,
    boxrule=2pt,
    leftrule=5pt,
    breakable
}

\newtcolorbox{warningbox}{
    colback=red!10,
    colframe=red!50,
    title=Warning: Common Pitfall,
    fonttitle=\bfseries,
    arc=3pt,
    boxrule=2pt,
    leftrule=5pt,
    breakable
}

\newtcolorbox{tipbox}{
    colback=green!10,
    colframe=green!50,
    title=Pro Tip,
    fonttitle=\bfseries,
    arc=3pt,
    boxrule=2pt,
    leftrule=5pt,
    breakable
}

\begin{document}

\title{AMC12B 2017 Problem 24: Strategic Problem Analysis}
\author{Competition Math Framework}
\date{\today}
\maketitle

\section*{Problem Statement}

Quadrilateral $ABCD$ has right angles at $B$ and $C$, $\triangle ABC \sim \triangle BCD$, and $AB > BC$. There is a point $E$ in the interior of $ABCD$ such that $\triangle ABC \sim \triangle CEB$ and the area of $\triangle AED$ is $17$ times the area of $\triangle CEB$. What is $\tfrac{AB}{BC}$?

\textbf{Answer Choices:} 
\begin{itemize}[nosep]
    \item (A) $1+\sqrt{2}$ 
    \item (B) $2 + \sqrt{2}$ 
    \item (C) $\sqrt{17}$ 
    \item (D) $2 + \sqrt{5}$ 
    \item (E) $1 + 2\sqrt{3}$
\end{itemize}

\begin{problemconfigbox}
\subsection*{A. Problem Classification}
\begin{itemize}[leftmargin=*]
    \item \textbf{Problem Type}: To Find (determine ratio $\frac{AB}{BC}$)
    \item \textbf{Difficulty Band}: Late-band (Problem 24 out of 25)
    \item \textbf{Competition Context}: 2017 AMC 12B
    \item \textbf{Mathematical Domain}: Geometry with algebraic computation
    \item \textbf{Sub-domain}: Similar triangles, quadrilaterals, area relationships
    \item \textbf{Time Allocation}: 8-12 minutes (second-hardest problem)
\end{itemize}

\subsection*{B. Problem Structure Identification}
\begin{itemize}[leftmargin=*]
    \item \textbf{Given Information}:
    \begin{itemize}
        \item Quadrilateral $ABCD$ with right angles at $B$ and $C$
        \item $\triangle ABC \sim \triangle BCD$ (two similar triangles in quadrilateral)
        \item $AB > BC$ (orientation constraint)
        \item Point $E$ in interior such that $\triangle ABC \sim \triangle CEB$
        \item $[AED] = 17 \cdot [CEB]$ (area relationship)
    \end{itemize}
    
    \item \textbf{Constraints}: 
    \begin{itemize}
        \item Three similarity conditions create strong geometric structure
        \item Right angles at $B$ and $C$ determine shape
        \item Area ratio of 17 provides algebraic constraint
        \item $E$ must be interior to $ABCD$
    \end{itemize}
    
    \item \textbf{Objective}: Find the ratio $\frac{AB}{BC}$
    
    \item \textbf{Hidden Assumptions}:
    \begin{itemize}
        \item From $\triangle ABC \sim \triangle BCD$, sides are in proportion
        \item If $BC = x$, then $AB = x^2$ and $CD = 1$ (by similarity)
        \item The similarity conditions determine $E$'s position uniquely
        \item Point $E$ creates additional similar triangles
    \end{itemize}
    
    \item \textbf{Answer Choice Leverage}: 
    \begin{itemize}
        \item All answers involve surds, suggesting algebraic/quartic equation
        \item Answer (C) $\sqrt{17}$ directly relates to the given area ratio
        \item Answer (D) $2 + \sqrt{5}$ suggests solving $x^4 - 18x^2 + 1 = 0$
    \end{itemize}
\end{itemize}

\subsection*{C. Strategic Orientation}
\begin{itemize}[leftmargin=*]
    \item \textbf{Hypothesis}: Set $BC = x$ as variable; use similarity to express all sides
    \item \textbf{Conclusion}: Solve for $x = \frac{AB}{BC}$ using area constraint
    \item \textbf{Notation}: 
    \begin{itemize}
        \item Let $CD = 1$, $BC = x$, $AB = x^2$ (from similarity scaling)
        \item Let $[T]$ denote area of triangle $T$
    \end{itemize}
    \item \textbf{Intuitive Approach}: Use coordinate system or similar triangle ratios
    \item \textbf{Cross-Domain Potential}: Could use coordinates, but pure geometry cleaner
\end{itemize}
\end{problemconfigbox}

\begin{strategicbox}
\subsection*{A. Psychological Strategies}
\begin{itemize}[leftmargin=*]
    \item \textbf{Mental Toughness}: Problem 24 - expect complex similarity chains and careful algebra
    \item \textbf{Creativity Cultivation}: Look for ways to express areas using similar triangle ratios
    \item \textbf{Iterative Refinement}: Start with simple cases, build up complexity systematically
\end{itemize}

\subsection*{B. Investigation Strategies}
\begin{itemize}[leftmargin=*]
    \item \textbf{Startup Orientation}: Draw careful diagram showing all right angles and similar triangles
    \item \textbf{Get Your Hands Dirty}: Set $BC = x$ and compute other sides from similarity
    \item \textbf{Penultimate Step}: Express all areas in terms of $x$, then use $[AED] = 17[CEB]$
    \item \textbf{Wishful Thinking}: Wish that setting $CD = 1$ normalizes the problem nicely
    \item \textbf{Make It Easier}: Use similarity ratios to avoid coordinate calculations
    \item \textbf{Visualization}: Key insight comes from seeing how $E$ divides the quadrilateral
\end{itemize}

\subsection*{C. Late-Band Strategic Playbook}
\begin{itemize}[leftmargin=*]
    \item \textbf{Answer-Choice Leverage}: Answer (C) $\sqrt{17}$ suggests checking if $x^2 = 17$
    \item \textbf{Penultimate Step}: The quartic $x^4 - 18x^2 + 1 = 0$ is key
    \item \textbf{Make It Easier}: Normalize by setting $CD = 1$
    \item \textbf{Conjecture}: From area ratio 17, expect $x^2 = 9 + 4\sqrt{5}$
\end{itemize}

\subsection*{D. Argument Construction Strategy}
\begin{itemize}[leftmargin=*]
    \item \textbf{Direct Proof}: Use similarity ratios to express all areas, solve equation
    \item \textbf{Key Lemmas}:
    \begin{enumerate}
        \item From $\triangle ABC \sim \triangle BCD$: sides scale by factor $x$
        \item From $\triangle ABC \sim \triangle CEB$: can determine $E$'s position
        \item Similar triangles share altitude ratios
    \end{enumerate}
\end{itemize}
\end{strategicbox}

\begin{tacticalbox}
\subsection*{Core Mathematical Tactics}

\textbf{T1: Similar Triangle Scaling}
\begin{itemize}[leftmargin=*]
    \item From $\triangle ABC \sim \triangle BCD$ with ratio $k$:
    \begin{itemize}
        \item $\frac{BC}{CD} = \frac{AB}{BC}$, so $BC^2 = AB \cdot CD$
        \item If $CD = 1$ and $BC = x$, then $AB = x^2$
    \end{itemize}
    \item From $\triangle ABC \sim \triangle CEB$:
    \begin{itemize}
        \item Same similarity ratio determines $CE$ and $BE$
        \item $BD = \sqrt{BC^2 + CD^2} = \sqrt{x^2 + 1}$
    \end{itemize}
\end{itemize}

\textbf{T2: Area Computation via Similar Triangles}
\begin{itemize}[leftmargin=*]
    \item If $\triangle T_1 \sim \triangle T_2$ with ratio $k$, then $[T_1] = k^2 [T_2]$
    \item Use altitudes from similar triangles to compute areas
    \item For right triangles: $[\triangle] = \frac{1}{2} \cdot \text{leg}_1 \cdot \text{leg}_2$
\end{itemize}

\textbf{T3: Pythagorean Theorem}
\begin{itemize}[leftmargin=*]
    \item Right angles at $B$ and $C$ enable distance calculations
    \item $BD^2 = BC^2 + CD^2 = x^2 + 1$
    \item Use to find positions and distances
\end{itemize}

\textbf{T4: Quadratic Formula and Nested Radicals}
\begin{itemize}[leftmargin=*]
    \item Final equation: $x^4 - 18x^2 + 1 = 0$
    \item Substitute $y = x^2$: $y^2 - 18y + 1 = 0$
    \item $y = \frac{18 \pm \sqrt{324 - 4}}{2} = \frac{18 \pm \sqrt{320}}{2} = 9 \pm 4\sqrt{5}$
    \item Since $AB > BC$, we need $x > 1$, so $x^2 = 9 + 4\sqrt{5}$
    \item Recognize: $9 + 4\sqrt{5} = (2 + \sqrt{5})^2$
\end{itemize}
\end{tacticalbox}

\begin{techniquesbox}
\subsection*{High-Probability Conversion Techniques}

\textbf{TM30: Similarity Ratio Scaling}
\begin{itemize}[leftmargin=*]
    \item \textbf{When}: Nested similar triangles where dimensions scale proportionally
    \item \textbf{Application}: From $\triangle ABC \sim \triangle BCD$:
    \begin{align*}
        \frac{AB}{BC} &= \frac{BC}{CD} \\
        AB \cdot CD &= BC^2
    \end{align*}
    \item Normalize: Set $CD = 1$, $BC = x$, then $AB = x^2$
    \item This transforms the problem into single-variable form
\end{itemize}

\textbf{TM5: Coordinate Geometry (Alternative)}
\begin{itemize}[leftmargin=*]
    \item Place $B$ at origin, $C$ at $(x, 0)$, $A$ at $(0, x^2)$, $D$ at $(x, 1)$
    \item Use similarity to find $E = (x_E, y_E)$
    \item Compute areas using coordinate formulas
    \item More computational but systematic
\end{itemize}

\textbf{TM13: Area Ratios from Altitudes}
\begin{itemize}[leftmargin=*]
    \item Key: Triangles sharing same altitude have area ratio $= $ base ratio
    \item Drop perpendiculars from $E$ to sides
    \item Use similar triangle properties: $\triangle CEB \sim \triangle CFE \sim \triangle EFB$
    \item Express $BF$ and $CF$ in terms of $x$
\end{itemize}

\textbf{TM12: Completing the Square / Perfect Square Recognition}
\begin{itemize}[leftmargin=*]
    \item Recognize $9 + 4\sqrt{5} = 4 + 4\sqrt{5} + 5 = (2)^2 + 2 \cdot 2 \cdot \sqrt{5} + (\sqrt{5})^2$
    \item Therefore $x^2 = (2 + \sqrt{5})^2$, so $x = 2 + \sqrt{5}$
\end{itemize}
\end{techniquesbox}

\begin{solutionbox}
\subsection*{Complete Solution}

\textbf{Step 1: Set Up Coordinate System and Use Similarity}

Let $CD = 1$, $BC = x$, and (from $\triangle ABC \sim \triangle BCD$) $AB = x^2$.

By Pythagorean theorem: $BD = \sqrt{BC^2 + CD^2} = \sqrt{x^2 + 1}$.

\textbf{Step 2: Find Position and Properties of Point $E$}

From $\triangle ABC \sim \triangle BCD$, we have established $CD = 1$, $BC = x$, $AB = x^2$.

By Pythagorean theorem: $BD = \sqrt{BC^2 + CD^2} = \sqrt{x^2 + 1}$.

From $\triangle ABC \sim \triangle CEB$ with correspondence $A \leftrightarrow C$, $B \leftrightarrow E$, $C \leftrightarrow B$:
\begin{itemize}
    \item Since $\angle ABC = 90^\circ$, we have $\angle CEB = 90^\circ$ (corresponding angles)
    \item The similarity ratio relates corresponding sides
\end{itemize}

Using the similarity ratio:
\begin{align*}
    \frac{CE}{BC} &= \frac{BC}{BD} \implies CE = \frac{BC^2}{BD} = \frac{x^2}{\sqrt{x^2+1}} \\
    \frac{BE}{AB} &= \frac{BC}{BD} \implies BE = \frac{AB \cdot BC}{AB \cdot BD/BC} = \frac{BC}{BD} \cdot BC = \frac{x}{\sqrt{x^2+1}}
\end{align*}

\textbf{Verification:} Check that $\angle CEB = 90^\circ$ using Pythagorean theorem:
\[CB^2 = CE^2 + BE^2 = \frac{x^4}{x^2+1} + \frac{x^2}{x^2+1} = \frac{x^4 + x^2}{x^2+1} = \frac{x^2(x^2+1)}{x^2+1} = x^2 \checkmark\]

\textbf{Step 3: Compute Areas Using Altitude Method}

Let $F$ be the foot of the perpendicular from $E$ to $\overline{BC}$.

From the altitude-on-hypotenuse theorem in right triangle $CEB$ with right angle at $E$:
\begin{align*}
    CF &= \frac{CE^2}{CB} = \frac{(x^2/\sqrt{x^2+1})^2}{x} = \frac{x^4/(x^2+1)}{x} = \frac{x^3}{x^2+1} \\
    BF &= \frac{BE^2}{CB} = \frac{(x/\sqrt{x^2+1})^2}{x} = \frac{x^2/(x^2+1)}{x} = \frac{x}{x^2+1}
\end{align*}

Verify: $CF + BF = \frac{x^3}{x^2+1} + \frac{x}{x^2+1} = \frac{x^3 + x}{x^2+1} = \frac{x(x^2+1)}{x^2+1} = x = BC$ \checkmark

\textbf{Key Observation - Equal Areas:}

Due to the symmetric construction with nested similar triangles, we have:
\[[CED] = [BEA] = [CEB] = \frac{x^3}{2(x^2+1)}\]

This equality arises because:
\begin{itemize}
    \item $CF$ serves as the base of $\triangle CED$ with $CD = 1$ parallel to $EF$
    \item $BF$ serves as the base of $\triangle BEA$ with $AB = x^2$ parallel to $EF$  
    \item The altitude from $D$ to line through $E$ parallel to $BC$ equals the altitude from $A$ to the same line
    \item By similar triangle properties and the symmetric positioning of $E$, all three triangles have equal areas
\end{itemize}

Computing $[CEB]$:
\[[CEB] = \frac{1}{2} \cdot CE \cdot BE = \frac{1}{2} \cdot \frac{x^2}{\sqrt{x^2+1}} \cdot \frac{x}{\sqrt{x^2+1}} = \frac{x^3}{2(x^2+1)}\]

\textbf{Step 4: Apply Area Partition and Given Constraint}

The quadrilateral $ABCD$ (a trapezoid with parallel sides $AB$ and $CD$) has area:
\[[ABCD] = \frac{1}{2}(AB + CD) \cdot BC = \frac{1}{2}(x^2 + 1) \cdot x = \frac{x^3 + x}{2}\]

Point $E$ partitions the quadrilateral into four triangles:
\[[ABCD] = [AED] + [DEC] + [CEB] + [BEA]\]

From our equal-area observation: $[CED] = [BEA] = [CEB] = \frac{x^3}{2(x^2+1)}$

Given constraint: $[AED] = 17[CEB]$

Substituting into the partition equation:
\begin{align*}
    [ABCD] &= 17[CEB] + [CEB] + [CEB] + [CEB] = 20[CEB] \\
    \frac{x^3 + x}{2} &= 20 \cdot \frac{x^3}{2(x^2+1)} \\
    \frac{x^3 + x}{2} &= \frac{10x^3}{x^2+1}
\end{align*}

\textbf{Step 5: Solve for $x$}

\begin{align*}
    (x^3 + x)(x^2 + 1) &= 20x^3 \\
    x^5 + x^3 + x^3 + x &= 20x^3 \\
    x^5 + 2x^3 + x &= 20x^3 \\
    x^5 - 18x^3 + x &= 0 \\
    x(x^4 - 18x^2 + 1) &= 0
\end{align*}

Since $x > 0$: $x^4 - 18x^2 + 1 = 0$

Let $y = x^2$:
\begin{align*}
    y^2 - 18y + 1 &= 0 \\
    y &= \frac{18 \pm \sqrt{324 - 4}}{2} = \frac{18 \pm \sqrt{320}}{2} = \frac{18 \pm 8\sqrt{5}}{2} = 9 \pm 4\sqrt{5}
\end{align*}

Since $AB > BC$, we need $x > 1$, so $x^2 = 9 + 4\sqrt{5}$.

\textbf{Step 6: Simplify the Radical}

Recognize that $9 + 4\sqrt{5}$ is a perfect square:
\begin{align*}
9 + 4\sqrt{5} &= 4 + 4\sqrt{5} + 5 \\
&= (2)^2 + 2(2)(\sqrt{5}) + (\sqrt{5})^2 \\
&= (2 + \sqrt{5})^2
\end{align*}

Therefore: $x = 2 + \sqrt{5}$

\[\boxed{\textbf{(D) } 2 + \sqrt{5}}\]

\textbf{Alternative Solution Methods:}

\textbf{Method 2 (Weighted Average with Parallel Line):} Draw line $FG$ through $E$ with $F$ on $BC$ and $G$ on $AD$, where $FG \parallel AB$. By the trapezoid midsegment generalization, $FG = \frac{1+x^4}{1+x^2}$. The ratio $FE:EG = [\triangle CBE]:[\triangle ADE] = 1:17$ follows from triangles sharing perpendicular height to $FG$. Since $FE = \frac{x^2}{1+x^2}$, we get $FG = 18 \cdot FE$, leading to the same quartic equation.

\textbf{Method 3 (Direct Area Formula):} Set up with $AB = a$, $BC = b$. Use similarity to express $CD = \frac{b^2}{a}$, then compute all areas directly using base $\times$ height formulas. The constraint $\frac{[AED]}{[CEB]} = 17$ leads to $\frac{a^2}{b^2} + \frac{b^2}{a^2} = 18$, giving the same quartic in the ratio $\frac{a}{b}$.
\end{solutionbox}

\begin{keyinsightbox}
\textbf{Key Insight}: The triple similarity condition ($\triangle ABC \sim \triangle BCD \sim \triangle CEB$) creates a powerful constraint that determines the shape uniquely. By normalizing with $CD = 1$ and setting $BC = x$, we get $AB = x^2$ from the first similarity. 

The critical observation is that the symmetric construction forces $[CED] = [BEA] = [CEB]$ - these three triangles have equal areas due to the nested similar triangle structure. This dramatically simplifies the area equation.

The area constraint $[AED] = 17[CEB]$ then produces the quartic $x^4 - 18x^2 + 1 = 0$, whose solution $x^2 = 9 + 4\sqrt{5}$ is recognizable as the perfect square $(2 + \sqrt{5})^2$.

\textbf{Recognition Pattern:} When you see $a + b\sqrt{c}$ as an answer to $x^2$, always check if it equals $(p + q\sqrt{c})^2$ by expanding and comparing coefficients. This pattern appears frequently in AMC late-band problems.
\end{keyinsightbox}

\begin{warningbox}
\textbf{Warning}: The most common pitfalls in this problem:

\begin{enumerate}
    \item \textbf{Incorrect Similarity Correspondence}: When stating $\triangle ABC \sim \triangle CEB$, be careful about which vertices correspond. The right angle at $B$ in $\triangle ABC$ corresponds to the right angle at $E$ in $\triangle CEB$.
    
    \item \textbf{Skipping the Equal Areas Justification}: Do not simply assert that $[CED] = [BEA] = [CEB]$ without understanding why. This comes from the altitude projections and symmetric construction, not from the similarity alone.
    
    \item \textbf{Algebraic Errors in Area Computation}: When computing $[CEB] = \frac{1}{2} \cdot CE \cdot BE$, make sure to simplify $\frac{x^2}{\sqrt{x^2+1}} \cdot \frac{x}{\sqrt{x^2+1}} = \frac{x^3}{x^2+1}$ correctly.
    
    \item \textbf{Choosing Wrong Root}: From $x^2 = 9 \pm 4\sqrt{5}$, we must choose the positive root $9 + 4\sqrt{5}$ because $AB > BC$ requires $x > 1$.
\end{enumerate}
\end{warningbox}

\begin{tipbox}
\textbf{Pro Tip}: When you arrive at $x^2 = 9 + 4\sqrt{5}$, immediately check if this is a perfect square. The pattern $a + b\sqrt{c} = (p + q\sqrt{c})^2$ occurs frequently in AMC problems. Expand $(2 + \sqrt{5})^2 = 4 + 4\sqrt{5} + 5 = 9 + 4\sqrt{5}$ to verify.
\end{tipbox}

\begin{verificationbox}
\subsection*{Verification Strategy}

\textbf{Level 5: Thorough Verification}

\begin{enumerate}
    \item \textbf{Check dimensions}: $\frac{AB}{BC} = 2 + \sqrt{5} \approx 4.236$
    \begin{itemize}
        \item Makes sense since $AB > BC$ is given
        \item Reasonable magnitude for similar triangle construction
    \end{itemize}
    
    \item \textbf{Verify quartic solution}:
    \begin{itemize}
        \item $x = 2 + \sqrt{5}$, so $x^2 = 9 + 4\sqrt{5}$
        \item $x^4 = (9 + 4\sqrt{5})^2 = 81 + 72\sqrt{5} + 80 = 161 + 72\sqrt{5}$
        \item Check: $x^4 - 18x^2 + 1 = 161 + 72\sqrt{5} - 18(9 + 4\sqrt{5}) + 1$
        \item $= 161 + 72\sqrt{5} - 162 - 72\sqrt{5} + 1 = 0$ \checkmark
    \end{itemize}
    
    \item \textbf{Verify similarity ratios}:
    \begin{itemize}
        \item $\frac{AB}{BC} = x$ and $\frac{BC}{CD} = x$ \checkmark
        \item Consistent scaling: $AB \cdot CD = BC^2$ gives $x^2 \cdot 1 = x^2$ \checkmark
    \end{itemize}
    
    \item \textbf{Verify perfect square expansion}:
    \begin{itemize}
        \item $(2 + \sqrt{5})^2 = 4 + 4\sqrt{5} + 5 = 9 + 4\sqrt{5}$ \checkmark
        \item This confirms our simplification is correct
    \end{itemize}
    
    \item \textbf{Answer choice confirmation}:
    \begin{itemize}
        \item Answer (D) $2 + \sqrt{5}$ matches our solution \checkmark
        \item Check that other answers don't satisfy the quartic:
        \begin{itemize}
            \item (A) $1+\sqrt{2} \approx 2.414$: $(1+\sqrt{2})^2 = 3 + 2\sqrt{2} \neq 9 + 4\sqrt{5}$
            \item (B) $2+\sqrt{2} \approx 3.414$: $(2+\sqrt{2})^2 = 6 + 4\sqrt{2} \neq 9 + 4\sqrt{5}$
            \item (C) $\sqrt{17} \approx 4.123$: $(\sqrt{17})^2 = 17 \neq 9 + 4\sqrt{5}$
            \item (E) $1+2\sqrt{3} \approx 4.464$: $(1+2\sqrt{3})^2 = 13 + 4\sqrt{3} \neq 9 + 4\sqrt{5}$
        \end{itemize}
    \end{itemize}
    
    \item \textbf{Sanity check on equal areas}:
    \begin{itemize}
        \item With $x = 2 + \sqrt{5}$, verify numerically that $[CED] = [BEA] = [CEB]$
        \item This provides additional confidence in the geometric construction
    \end{itemize}
\end{enumerate}
\end{verificationbox}

\begin{competitionbox}
\subsection*{Competition Context and Timing}

\textbf{Time Management}
\begin{itemize}[leftmargin=*]
    \item \textbf{Target Time}: 8-12 minutes
    \item \textbf{Recognition Phase}: 2 minutes (identify triple similarity)
    \item \textbf{Setup Phase}: 3-4 minutes (normalize variables, express areas)
    \item \textbf{Calculation Phase}: 4-5 minutes (derive and solve quartic)
    \item \textbf{Verification Phase}: 1-2 minutes (check perfect square)
\end{itemize}

\textbf{Difficulty Assessment}
\begin{itemize}[leftmargin=*]
    \item One of the two hardest problems (24/25)
    \item Requires multiple advanced techniques:
    \begin{itemize}
        \item Similar triangle chain reasoning
        \item Area computations with constraints
        \item Quartic equation solving
        \item Perfect square recognition
    \end{itemize}
    \item Most students will not attempt this problem
    \item High risk/high reward - only for strongest geometry students
\end{itemize}

\textbf{Abandonment Criteria}
\begin{itemize}[leftmargin=*]
    \item If unable to set up similarity ratios after 3 minutes, skip
    \item If stuck on area computations after 6 minutes, guess and move on
    \item If quartic seems unsolvable, try answer choices
\end{itemize}

\textbf{Strategic Positioning}
\begin{itemize}[leftmargin=*]
    \item Attempt ONLY after completing Problems 1-20
    \item This problem distinguishes top 1\% scorers
    \item Perfect execution can secure AIME qualification
    \item Consider spending time on Problem 25 instead if geometry weak
\end{itemize}
\end{competitionbox}

\section*{Final Answer}
\[\boxed{2 + \sqrt{5}}\]

\section*{Solution Summary}

\textbf{Core Technique Chain:}
\begin{enumerate}
    \item \textbf{Normalization}: Set $CD = 1$, $BC = x$, then $AB = x^2$ from similarity
    \item \textbf{Similarity Application}: Use $\triangle ABC \sim \triangle CEB$ to find $CE$, $BE$
    \item \textbf{Equal Areas Recognition}: Identify $[CED] = [BEA] = [CEB]$ from symmetric construction
    \item \textbf{Area Constraint}: Apply $[AED] = 17[CEB]$ to get $20[CEB] = [ABCD]$
    \item \textbf{Algebraic Reduction}: Derive quartic $x^4 - 18x^2 + 1 = 0$
    \item \textbf{Perfect Square Recognition}: Simplify $x^2 = 9 + 4\sqrt{5} = (2+\sqrt{5})^2$
\end{enumerate}

\textbf{Key Success Factors:}
\begin{itemize}
    \item Recognizing the power of normalization early
    \item Understanding that nested similarities create equal-area regions
    \item Systematic algebraic manipulation without computational errors
    \item Pattern recognition for perfect squares with nested radicals
\end{itemize}

\section*{Confidence Assessment}
\begin{itemize}[leftmargin=*]
    \item \textbf{Confidence Level}: 98\% (elegant solution path, multiple verification methods confirm answer)
    
    \item \textbf{Verification Applied}: 
    \begin{itemize}
        \item Quartic equation verified by substitution
        \item Perfect square expansion verified
        \item All similarity ratios checked for consistency
        \item Answer choices tested against the quartic
        \item Dimensional analysis confirms $x > 1$
    \end{itemize}
    
    \item \textbf{Flag for Review}: No (very high confidence in both geometry and algebra)
    
    \item \textbf{Time Spent}: Approximately 10 minutes for complete solution
    \begin{itemize}
        \item 2 min: Problem analysis and normalization setup
        \item 3 min: Similarity calculations and area expressions
        \item 3 min: Algebraic manipulation to quartic
        \item 2 min: Perfect square recognition and verification
    \end{itemize}
    
    \item \textbf{Alternative Methods Available}: 
    \begin{itemize}
        \item Weighted average with parallel line (Solution 2 from AoPS)
        \item Coordinate geometry approach (Solution 4 from AoPS)
        \item Direct area computation (Solution 3 from AoPS)
    \end{itemize}
    
    \item \textbf{Competition Context}: This is one of the two hardest problems on the exam. Success on this problem typically indicates AIME qualification level mathematics. The solution requires mastery of:
    \begin{itemize}
        \item Similar triangle chains
        \item Strategic variable normalization
        \item Area computation techniques
        \item Quartic equation solving
        \item Nested radical simplification
    \end{itemize}
\end{itemize}

\textbf{Final Recommendation:} This problem is suitable for advanced students (AIME-level) who have strong geometry foundations. The solution presented here is the most elegant approach, but students should be aware of alternative methods for verification purposes.

\end{document}